% =====================================================================
%  LDA-T318 — Lab 3 briefing — Serve and measure
%  Thursday 12 November 2026
%  Build:  latexmk -xelatex main.tex
% =====================================================================
\documentclass[aspectratio=169,10pt]{beamer}

% Path has no parent-directory step: Overleaf resolves \input from the project root and
% will not read above it. Locally, the Makefile puts common/ on TEXINPUTS.
% =====================================================================
%  LDA-T318 — Large Language Models and Generative AI for NLP
%  Shared Beamer preamble for all lecture decks.
%  Aarne Talman, University of Helsinki, 2026.
%
%  Engine:  XeLaTeX or LuaLaTeX preferred (gives the Fira font set).
%           pdfLaTeX also works; Metropolis falls back to Latin Modern.
%  Build:   latexmk -xelatex main.tex      (from the deck's own directory)
%
%  Every deck does:
%      \documentclass[aspectratio=169,10pt]{beamer}
%      \input{../common/preamble.tex}
% =====================================================================

% ---------------------------------------------------------------- theme
\usetheme{metropolis}
\metroset{
  titleformat=smallcaps,
  numbering=fraction,
  progressbar=frametitle,
  block=fill,
  sectionpage=progressbar,
  subsectionpage=none,
}

% ---------------------------------------------------------------- packages
% inputenc/fontenc are required under pdfLaTeX and actively harmful under
% XeLaTeX/LuaLaTeX, where Metropolis loads fontspec to get Fira. Guard them.
\usepackage{iftex}
\ifPDFTeX
  \usepackage[utf8]{inputenc}
  \usepackage[T1]{fontenc}
  \usepackage[sfdefault]{FiraSans}
  \usepackage{FiraMono}
\else
  % Metropolis only gets Fira if fontspec can actually find it. TeX Live ships
  % the OTFs but does NOT register them with fontconfig, so a plain
  % \setsansfont{Fira Sans} silently falls back to Latin Modern and the deck
  % stops looking like Metropolis. Load them by filename instead.
  \usepackage{fontspec}
  \setsansfont{FiraSans}[
    Extension      = .otf,
    UprightFont    = *-Book,
    ItalicFont     = *-BookItalic,
    BoldFont       = *-Medium,
    BoldItalicFont = *-MediumItalic]
  \setmonofont{FiraMono}[
    Extension   = .otf,
    UprightFont = *-Regular,
    BoldFont    = *-Medium,
    Scale       = 0.92]
  \renewcommand{\familydefault}{\sfdefault}
\fi
\usepackage[english]{babel}

\usepackage{amsmath,amssymb}
\usepackage{booktabs}
\usepackage{array}
\usepackage{multicol}
\usepackage{xcolor}
\usepackage{graphicx}
\usepackage{tikz}
\usepackage{pgfplots}
\usepackage{listings}
\usepackage{appendixnumberbeamer}

\pgfplotsset{compat=1.17}
\usetikzlibrary{positioning,arrows.meta,fit,backgrounds,calc,matrix,shapes.geometric,decorations.pathreplacing}

% =====================================================================
%  University of Helsinki palette
%  Primary brand colour is UH Hray/black with the "HY blue" accent.
%  Secondary accents chosen for projector contrast, not for print.
% =====================================================================
\definecolor{UHBlue}{HTML}{0C5A9E}   % primary accent
\definecolor{UHDark}{HTML}{15202B}   % body text / dark backgrounds
\definecolor{UHGrey}{HTML}{5A6672}   % secondary text
\definecolor{UHLight}{HTML}{EEF2F6}  % block backgrounds
\definecolor{UHOrange}{HTML}{D1601A} % alerts, gotchas
\definecolor{UHGreen}{HTML}{1F7A5C}  % examples, "this works"
\definecolor{UHRed}{HTML}{A62B2B}    % failure modes
\definecolor{UHSand}{HTML}{C8B88A}   % neutral third accent

\setbeamercolor{normal text}{fg=UHDark, bg=white}
\setbeamercolor{alerted text}{fg=UHOrange}
\setbeamercolor{example text}{fg=UHGreen}
\setbeamercolor{frametitle}{fg=white, bg=UHBlue}
\setbeamercolor{title separator}{fg=UHBlue}
\setbeamercolor{progress bar}{fg=UHBlue}
\setbeamercolor{progress bar in head/foot}{fg=UHBlue}
\setbeamercolor{progress bar in section page}{fg=UHBlue}
\setbeamercolor{structure}{fg=UHBlue}
\setbeamercolor{palette primary}{fg=white, bg=UHBlue}   % section pages, standout frames
\setbeamercolor{block title}{fg=white, bg=UHBlue}
\setbeamercolor{block body}{fg=UHDark, bg=UHLight}
\setbeamercolor{block title alerted}{fg=white, bg=UHOrange}
\setbeamercolor{block body alerted}{fg=UHDark, bg=UHLight}
\setbeamercolor{block title example}{fg=white, bg=UHGreen}
\setbeamercolor{block body example}{fg=UHDark, bg=UHLight}
\setbeamercolor{footnote}{fg=UHGrey}

\setbeamerfont{frametitle}{size=\large, series=\bfseries}
\setbeamerfont{footnote}{size=\tiny}
\setbeamerfont{caption}{size=\scriptsize}

% Frame title bar.
%
% Do NOT reimplement the frametitle template here. Metropolis's own template
% already spans \paperwidth (so `bg=UHBlue' above gives the full-width bar),
% already struts two-line titles without collision, and already draws the
% frametitle progress bar. Reimplementing it means touching internals whose
% names contain `@', which is NOT a letter in a .tex file -- `\@empty' parses
% as `\@' followed by the word "empty" and the build breaks.
%
% To change the bar's height, adjust Metropolis's own padding length instead,
% inside \makeatletter so that `@' is a letter:
%
%   \makeatletter
%   \setlength{\metropolis@frametitle@padding}{2.4ex}  % default 2.2ex
%   \makeatother
%
\makeatletter
\setlength{\metropolis@frametitle@padding}{1.2ex}
\makeatother

% No deck uses \framesubtitle; if one ever does, Metropolis ignores it and the
% subtitle should go in the frame body instead.

% =====================================================================
%  Course metadata — decks override \lecturename and \lecturenumber
% =====================================================================
\newcommand{\coursecode}{LDA-T318}
\newcommand{\coursename}{LLMs and Generative AI for NLP}
\newcommand{\courseteacher}{Aarne Talman}
\newcommand{\courseinstitution}{University of Helsinki}
\providecommand{\lecturename}{}
\providecommand{\lecturedate}{}

% Footline: course code on the left, lecture name centred, n/N on the right.
\setbeamertemplate{frame footer}{\coursecode\quad\textbar\quad\lecturename}

% =====================================================================
%  Semantic macros — use these instead of ad-hoc formatting, so that a
%  change of house style is a single edit here.
% =====================================================================

% --- Sourcing. Every quantitative claim on a slide should carry one. ---
\newcommand{\source}[1]{%
  \vfill{\scriptsize\color{UHGrey}\raggedright Source: #1\par}%
}
\newcommand{\sources}[1]{%
  \vfill{\scriptsize\color{UHGrey}\raggedright Sources: #1\par}%
}
% Inline citation chip, e.g. \cit{Raschka 2026a}
\newcommand{\cit}[1]{{\scriptsize\color{UHGrey}[#1]}}

% --- Counter-intuitive point. Three per lecture; these are exam material. ---
\newenvironment{gotcha}[1][Gotcha]{%
  \begin{alertblock}{#1}%
}{%
  \end{alertblock}%
}

% --- Speculative / unconfirmed claim. Research literacy matters here. ---
\newcommand{\rumour}{{\color{UHOrange}\textbf{[unconfirmed]}}\ }

% --- Key takeaway, for the closing slide of each section. ---
\newenvironment{takeaway}{%
  \begin{block}{Take away}%
}{%
  \end{block}%
}

% --- Live demo marker. ---
\newcommand{\demo}[1]{{\color{UHGreen}\textbf{Demo:}} #1}

% --- Emphasis consistent with the palette. ---
\newcommand{\key}[1]{{\color{UHBlue}\textbf{#1}}}
\newcommand{\warn}[1]{{\color{UHOrange}\textbf{#1}}}
\newcommand{\good}[1]{{\color{UHGreen}\textbf{#1}}}
\newcommand{\bad}[1]{{\color{UHRed}\textbf{#1}}}

% --- Units and quantities used constantly across the decks. ---
\newcommand{\unit}[1]{\,\mathrm{#1}}
\newcommand{\flops}{\mathrm{FLOPs}}
% NB: no \toks macro here --- \toks is a TeX primitive (a token register), so
% \newcommand{\toks} is a hard "already defined" error. No deck used it.
\newcommand{\KV}{\mathrm{KV}}
\newcommand{\softmax}{\operatorname{softmax}}

% =====================================================================
%  Code listings — used for config snippets, Slurm scripts, chat templates.
%  Frames containing these MUST be declared [fragile].
% =====================================================================
\lstdefinestyle{uh}{
  basicstyle=\ttfamily\scriptsize\color{UHDark},
  keywordstyle=\color{UHBlue}\bfseries,
  commentstyle=\color{UHGrey}\itshape,
  stringstyle=\color{UHGreen},
  numberstyle=\tiny\color{UHGrey},
  backgroundcolor=\color{UHLight},
  frame=none,
  xleftmargin=0.6em,
  framexleftmargin=0.6em,
  breaklines=true,
  breakatwhitespace=true,
  showstringspaces=false,
  columns=fullflexible,
  keepspaces=true,
  upquote=true,
  literate={-}{{-}}1 {_}{{\_}}1,
}
\lstset{style=uh}

\lstdefinelanguage{slurm}{
  morekeywords={SBATCH,module,srun,source,export,set,echo},
  morecomment=[l]{\#},
  sensitive=true,
}
\lstdefinelanguage{jinja}{
  morekeywords={for,endfor,if,endif,set,in},
  morecomment=[s]{\{\#}{\#\}},
  sensitive=true,
}

% =====================================================================
%  TikZ styles shared across decks (block diagrams, token grids).
% =====================================================================
\tikzset{
  mod/.style     = {draw=UHBlue, fill=UHLight, rounded corners=2pt,
                    minimum height=7mm, minimum width=22mm,
                    font=\scriptsize, align=center},
  modhi/.style   = {mod, draw=UHOrange, fill=UHOrange!12},
  modmute/.style = {mod, draw=UHGrey, fill=white, text=UHGrey},
  flow/.style    = {-{Latex[length=2mm]}, draw=UHGrey, thick},
  flowhi/.style  = {flow, draw=UHOrange},
  resid/.style   = {-{Latex[length=2mm]}, draw=UHBlue, thick, dashed},
  tokcell/.style = {draw=UHGrey!50, minimum size=4.2mm, font=\tiny, inner sep=0pt},
  tokon/.style   = {tokcell, fill=UHBlue!35},
  tokoff/.style  = {tokcell, fill=UHLight},
  tokmask/.style = {tokcell, fill=UHGrey!15, text=UHGrey!60},
  lbl/.style     = {font=\scriptsize\color{UHGrey}},
}

% =====================================================================
%  Presenter notes.
%
%  Notes are written on every slide. By default they are NOT shown, so
%  `latexmk -xelatex main.tex` produces a clean deck for projection.
%
%  To present with notes on a second screen (pdfpc, Skim, Impressive):
%      uncomment the line below, rebuild to main-notes.pdf
%  To print a notes handout:
%      \setbeameroption{show only notes}
% =====================================================================
% \setbeameroption{show notes on second screen=right}

\setbeamertemplate{note page}{%
  {\scriptsize\color{UHGrey}\coursecode\ \textbar\ \lecturename\par}%
  \vskip0.3em\hrule\vskip0.6em%
  {\small\insertnote}%
}

% ---------------------------------------------------------------- misc
\setbeamertemplate{caption}[numbered]
\setbeamertemplate{itemize items}[circle]
\setlength{\parskip}{0.1em}

% ---------------------------------------------------------------- density
% These decks are text-dense and a 16:9 frame is only 90mm tall, so the
% default vertical rhythm costs several lines per slide. Tightened globally
% here rather than slide by slide.
\setbeamersize{text margin left=0.7cm, text margin right=0.7cm}
\makeatletter
\def\@listi{\leftmargin\leftmargini
  \topsep 1.5\p@ \@plus1\p@ \@minus1\p@
  \parsep 0\p@
  \itemsep 1.5\p@ \@plus1\p@ \@minus1\p@}
\let\@listI\@listi
\makeatother

% The title page is a stack of short paragraphs, so a non-zero \parskip
% pushes it past the bottom of the frame. Neutralise it there only.
\addtobeamertemplate{title page}{\setlength{\parskip}{0pt}}{}

% Tighter lists — the decks are dense.
\setbeamertemplate{itemize/enumerate body begin}{\small}
\setbeamertemplate{itemize/enumerate subbody begin}{\footnotesize}

% Table helpers
\newcolumntype{L}[1]{>{\raggedright\arraybackslash}p{#1}}
\newcolumntype{C}[1]{>{\centering\arraybackslash}p{#1}}
\newcolumntype{R}[1]{>{\raggedleft\arraybackslash}p{#1}}

% Allow slightly overfull boxes without a warning storm
\hfuzz=2pt

% =====================================================================
%  Shared course-map frame.
%  Every lecture opens with this, with the current week highlighted:
%      \coursemapframe{1}   ... \coursemapframe{4}
%  Pass 0 to highlight nothing.
% =====================================================================
% \cmrow{current week}{this row's week}{dates}{Tuesday}{Thursday}{teacher}
% Highlights the row in UH blue bold when the two week numbers match.
% Deliberately avoids \rowcolor so that colortbl is not needed.
% The & and \\ must NOT sit inside the \ifnum branches: TeX scans alignment
% cells while skipping the false branch, so an & there breaks the conditional
% ("Incomplete \ifnum"). Keep the branches down to a macro name and let the
% separators live in the macro bodies, which expand after the test resolves.
\newcommand{\cmhi}[5]{%
  \textcolor{UHBlue}{\textbf{#1}} & \textcolor{UHBlue}{\textbf{#2}} &
  \textcolor{UHBlue}{\textbf{#3}} & \textcolor{UHBlue}{\textbf{#4}} &
  \textcolor{UHBlue}{\textbf{#5}}\\%
}
\newcommand{\cmplain}[5]{#1 & #2 & #3 & #4 & #5\\}
\newcommand{\cmrow}[6]{%
  \ifnum#1=#2\relax\expandafter\cmhi\else\expandafter\cmplain\fi
  {#2}{#3}{#4}{#5}{#6}%
}
\newcommand{\coursemapframe}[1]{%
\begin{frame}{Course overview}
\centering
\small
\setlength{\parskip}{0pt}
\renewcommand{\arraystretch}{1.08}
\begin{tabular}{@{}c L{1.7cm} L{3.1cm} L{3.7cm} L{1.4cm}@{}}
\toprule
\textbf{Wk} & \textbf{Dates} & \textbf{Tuesday} & \textbf{Thursday} & \textbf{Teacher}\\
\midrule
\cmrow{#1}{1}{27 / 29 Oct}{Language models}{Pretraining}{Aarne}
\cmrow{#1}{2}{3 / 5 Nov}{Post-training}{Lab 2: tune vs prompt}{Aarne}
\cmrow{#1}{3}{10 / 12 Nov}{Inference}{Lab 3: serve}{Aarne}
\cmrow{#1}{4}{17 / 19 Nov}{Evaluation}{Lab 4: evaluate}{Jussi}
\cmrow{#1}{5}{24 / 26 Nov}{RAG}{Lab 5: RAG}{Dmitry}
\cmrow{#1}{6}{1 / 3 Dec}{Agents}{Lab 6: agent}{Dmitry}
\cmrow{#1}{7}{8 / 10 Dec}{Course recap}{Special topics (TBC)}{All}
\bottomrule
\end{tabular}

\vspace{1.2em}
{\footnotesize Reusing a task and model across labs is recommended, not required.\\
There are no intermediate submissions; the final report and code repository
determine the grade.}
\note{Fifteen seconds, not more. The point of this slide is orientation, not
content: show them the shape of the course and where today sits inside it.
The final week is deliberately left flexible for a recap and special topics.}
\end{frame}%
}


\renewcommand{\lecturename}{Lab 3 — Serve and measure}
\renewcommand{\lecturedate}{12 November 2026}

\title{Lab 3: Serve and measure}
\subtitle{Measuring what your model costs to run}
\author{Aarne Talman}
\institute{\coursecode\ \textbar\ \coursename\ \textbar\ \courseinstitution}
\date{12 November 2026}

\begin{document}

% =====================================================================
\begin{frame}[plain,shrink=8]
  \titlepage
  \note{Fifteen minutes of briefing, then the room works and I circulate.
  Say that out loud in the first thirty seconds so nobody settles in for a
  lecture.

  The frame to hold: on Tuesday we did the theory of serving, and today every
  number on those slides becomes a number they measured themselves on their
  own model. Nothing gets trained today; the session is all measurement.

  Target: briefing done by minute 10, every server in the room answering
  requests by minute 15.}
\end{frame}

\coursemapframe{3}

% =====================================================================
\section{The question}
% =====================================================================

\begin{frame}{The question today}
  Today you put a language model behind a server and measure what it costs to
  run. You may use the adapter from Lab 2 or another compatible model.
  \vspace{0.4em}

  \begin{block}{Today}
    \textbf{What does it cost to run, and what do you give up to make it
    cheaper?}
  \end{block}
  \vspace{0.4em}

  \begin{columns}[T]
    \begin{column}{0.5\textwidth}
      \textbf{You will produce}
      \begin{itemize}
        \item one \key{speed-versus-quality table}
        \item a \key{deployment recommendation}, 3--5 sentences
        \item a \key{one-command restart script} you can reuse in later labs
      \end{itemize}
    \end{column}
    \begin{column}{0.5\textwidth}
      \textbf{You will not}
      \begin{itemize}
        \item train anything
        \item change your prompt
        \item change your held-out set
      \end{itemize}
      \vspace{0.3em}
      Every number today is a measurement on a model you already have.
    \end{column}
  \end{columns}
  \vfill
  \small In your final report, the deployment recommendation must be supported
  by the table. Start with \texttt{lab3-example/README.md}.
  \note{Thirty seconds. The sentence to land: cost and quality are a single
  trade-off, and today they finally sit in the same table.

  Explain what a useful recommendation includes: which configuration to
  deploy, under what latency target, what would change at ten times the
  traffic, and what FP8 cost in quality. Three to five sentences, backed by
  the rows of the table.

  While circulating: the students who finish early are the ones who skipped the
  warm-up. Check their numbers before you let them write the recommendation.}
\end{frame}

\begin{frame}{Material for the final report}
  \centering
  \scriptsize
  \begin{tabular}{@{}L{2.3cm} C{1.5cm} C{1.5cm} C{1.9cm} C{1.2cm} C{1.2cm}@{}}
    \toprule
    \textbf{Configuration} & \textbf{TTFT ms, p50 / p95} & \textbf{TPOT ms, p50 / p95} &
    \textbf{Output tok/s @ 1 / 8 / 32} & \textbf{Task score} & \textbf{Invalid \%}\\
    \midrule
    bf16 + LoRA adapter & \color{UHGrey}-- & \color{UHGrey}-- & \color{UHGrey}-- &
      \color{UHGrey}-- & \color{UHGrey}--\\[0.3em]
    FP8 + LoRA adapter & \color{UHGrey}-- & \color{UHGrey}-- & \color{UHGrey}-- &
      \color{UHGrey}-- & \color{UHGrey}--\\[0.3em]
    bf16 + JSON schema & \color{UHGrey}-- & \color{UHGrey}-- & \color{UHGrey}-- &
      \color{UHGrey}-- & \color{UHGrey}--\\
    \bottomrule
  \end{tabular}

  \vspace{0.6em}
  \raggedright\scriptsize
  \textbf{TTFT}: how long until the \emph{first} token appears.
  \textbf{TPOT}: milliseconds per token after that.
  \textbf{p50 / p95}: the median request, and the slowest one in twenty.
  \textbf{@ 1 / 8 / 32}: with that many requests in flight at the same instant.
  \textbf{Invalid \%}: outputs your parser could not use.
  \vspace{0.4em}

  \small
  Same model, same adapter, same held-out set, same prompts, same decoding
  parameters in every row. \key{One variable changes per row.}
  Task score comes from the \textbf{same held-out scorer used in week 2},
  unchanged.
  \note{Put this on the screen and leave it there while you talk. It is the
  shape of the whole session: three rows, and each block of the timeline fills
  one of them. The three definitions along the bottom are first use, so read
  them out; the whole second half of the briefing assumes them. Say that
  ``number of requests behind every latency cell'' is not optional. A
  percentile computed from five requests tells you nothing.

  Say explicitly that the third row only needs the invalid-output and score
  columns if time is short; speed under constrained decoding is a bonus.

  The mistake to watch for while circulating: students who change two things
  between rows, typically precision and max-model-len, because the FP8 server
  would not start at the old setting. Then the comparison is worthless. Tell
  them to go back and re-run the bf16 row at the new setting instead.}
\end{frame}

\begin{frame}{Ninety minutes}
  \centering
  \vspace{0.6em}
  \begin{tikzpicture}[xscale=0.122]
    % --- blocks -----------------------------------------------------
    \fill[UHOrange!15] (0,0)  rectangle (10,1.1);
    \draw[UHOrange]    (0,0)  rectangle (10,1.1);
    \node[font=\tiny, align=center, text width=1.0cm] at (5,0.55)
      {start the server};

    \fill[UHLight] (10,0) rectangle (35,1.1);
    \draw[UHBlue]  (10,0) rectangle (35,1.1);
    \node[font=\scriptsize, align=center, text width=2.8cm] at (22.5,0.55)
      {speed: TTFT, TPOT, throughput at 1 / 8 / 32};

    \fill[UHLight] (35,0) rectangle (60,1.1);
    \draw[UHBlue]  (35,0) rectangle (60,1.1);
    \node[font=\scriptsize, align=center, text width=2.8cm] at (47.5,0.55)
      {FP8: re-time and re-score};

    \fill[UHLight] (60,0) rectangle (80,1.1);
    \draw[UHBlue]  (60,0) rectangle (80,1.1);
    \node[font=\scriptsize, align=center, text width=2.2cm] at (70,0.55)
      {JSON schema: invalid-output rate};

    \fill[UHGreen!12] (80,0) rectangle (90,1.1);
    \draw[UHGreen]    (80,0) rectangle (90,1.1);
    \node[font=\tiny, align=center, text width=1.0cm] at (85,0.55)
      {save and write up};

    % --- axis -------------------------------------------------------
    \draw[UHGrey] (0,-0.1) -- (90,-0.1);
    \foreach \x in {0,10,35,60,80,90} {
      \draw[UHGrey] (\x,0) -- (\x,-0.25);
      \node[font=\tiny, text=UHGrey] at (\x,-0.5) {\x};
    }
    \node[font=\tiny, text=UHGrey, anchor=west] at (91,-0.5) {min};
  \end{tikzpicture}
  \vspace{0.9em}

  \raggedright\small
  If the server is not up by \textbf{minute 15}, stop debugging alone and put
  your hand up. Everything after that block depends on it.
  \note{Read the timeline out loud once. It is the only schedule they get and
  it stays on the wall at the end of the deck.

  The honest warning: the first block is the one that overruns. Weight
  download, a stale kernel holding the GPU, a port already in use. That is why
  the troubleshooting appendix exists. Know it well enough to diagnose from the
  doorway.

  Circulating plan: minutes 0--15 is triage, nothing else. Walk the room once
  and look only at whether a server is answering. From minute 15 switch to
  looking at measurement method.}
\end{frame}

% =====================================================================
\section{Serving}
% =====================================================================

% Maintained commands and client code live in lab3-example.
\iffalse
\begin{frame}[fragile,shrink=12]{Start the server}
\begin{lstlisting}[language=slurm]
# Run this inside a gpuinteractive session (one GH200). NOT on a login node.
module purge                       # start from an empty module environment
module load python-vllm            # vLLM has its own module, loaded in week 3 only

export HF_HOME=/scratch/project_2021893/hf  # shared weight cache

vllm serve Qwen/Qwen3.5-4B-Instruct \
  --served-model-name base \
  --enable-lora \
  --lora-modules task=/scratch/project_2021893/lora/w2-adapter \
  --max-model-len 4096 \
  --gpu-memory-utilization 0.85 \
  --port 8000
\end{lstlisting}
  \vspace{0.2em}
  \small
  \textbf{Success looks like} the line \texttt{Application startup complete}
  (it can take a few minutes the first time, while the weights download). Check
  from a second terminal: \texttt{curl localhost:8000/v1/models} must list
  \texttt{base} \emph{and} \texttt{task}.
  \vspace{0.3em}

  \scriptsize
  \texttt{--lora-modules task=<path>} loads your week 2 adapter and names it
  \texttt{task}; that name goes in every request. \quad
  \texttt{--max-model-len} caps prompt plus answer, in tokens. \quad
  \texttt{--gpu-memory-utilization} is a \key{reservation}: vLLM claims
  85\,\% of the card up front, and whatever is left after the weights becomes
  KV cache, the per-request memory that decides how many requests fit at once.
  Tuesday's lecture was about that preallocation.
  \note{Walk four lines only: the module, the adapter, max-model-len and
  gpu-memory-utilization. Everything else they can read.

  The adapter line is the one people get wrong. The name to the left of the
  equals sign is what they will put in the model field of every request. Say it
  now and say it again on the next slide.

  The memory line connects straight to Tuesday: the weights are about 8 GB in
  bf16, so the other seventy-odd gigabytes are KV cache, which is what caps
  concurrency. If anyone asks why not 0.95: leave headroom for the eval script
  and the notebook kernel.

  Watch for: people serving the base model when their adapter was trained on
  the instruct model. Ask them which checkpoint they fine-tuned.}
\end{frame}

\begin{frame}[fragile,shrink=15]{Call it from the notebook}
\begin{lstlisting}[language=Python]
from openai import OpenAI
import time

client = OpenAI(base_url="http://localhost:8000/v1", api_key="EMPTY")  # key unused

def one_request(prompt, n_out=128):
    t0 = time.perf_counter()
    stream = client.chat.completions.create(
        model="task",                      # the adapter name, not the repo id
        messages=[{"role": "user", "content": prompt}],
        temperature=0.0, max_tokens=n_out, stream=True,
        extra_body={"ignore_eos": True},   # exactly n_out tokens, always
    )
    ttft, n = None, 0
    for chunk in stream:
        if not chunk.choices[0].delta.content:
            continue                       # skip the role-only first chunk
        if ttft is None: ttft = time.perf_counter() - t0
        n += 1
    return ttft, (time.perf_counter() - t0 - ttft) / max(n - 1, 1)  # TTFT, TPOT
\end{lstlisting}
  \small
  \key{You must stream} to measure TTFT at all. A non-streaming call hands you
  the whole answer at once, so there is no ``first token'' to time.
  \textbf{Success looks like} a TTFT of a few tens of milliseconds and a TPOT
  well under it; if TTFT comes back equal to the total, streaming is off.
  \note{This function is the whole measurement harness. Everything later today
  is this function called in a loop, or called from several threads at once.

  Two things to point at: model equals the adapter name, and stream equals
  True. Without streaming there is no TTFT, and about a third of the room will
  try to compute it from the total anyway. The content check inside the loop
  matters too: the first chunk carries only the role, so counting it puts TTFT
  one chunk too early.

  Note the api-key placeholder. The server does not check it, but the OpenAI
  client refuses to start without one. This is also exactly the client Dmitry's
  labs will use in weeks 5 and 6.

  Watch for: a notebook kernel still holding GPU memory from week 2, so the
  server never started and every call raises a connection error. The fix is to
  restart the kernel rather than the server.}
\end{frame}

% =====================================================================
\section{Measuring}
% =====================================================================

\fi
\begin{frame}{Measure honestly, 1: warm up, then fix the length}
  \begin{columns}[T]
    \begin{column}{0.52\textwidth}
      \textbf{Before you time anything}
      \begin{itemize}
        \item send $\approx$ 20 throwaway requests and \key{discard them}
        \item then start recording
      \end{itemize}
      \vspace{0.5em}
      \textbf{Fix everything except the one variable}
      \begin{itemize}
        \item \texttt{max\_tokens=128} and \texttt{ignore\_eos=True}, so that
          every request emits exactly 128 tokens
        \item the same prompt set, in the same order
        \item \texttt{temperature=0.0} everywhere
      \end{itemize}
      \vspace{0.4em}
      Without a fixed output length you are measuring \emph{how verbose the
      model is} rather than how fast the server is. FP8 changes verbosity.
    \end{column}
    \begin{column}{0.48\textwidth}
      \begin{gotcha}[The first request is always slow]
        One-off work all lands on request number one: the GPU compiles and
        tunes its kernels, the weights are read in for real, and nothing is
        cached yet.
        \vspace{0.3em}
        Without a warm-up you are measuring \textbf{startup}.
      \end{gotcha}
      \vspace{0.3em}
      \small
      The opposite trap: the server remembers the work done for a prompt it has
      already seen (``prefix caching''), so the \emph{second} identical prompt
      looks artificially fast. Vary your prompts, or serve with
      \texttt{--no-enable-prefix-caching} while benchmarking.
    \end{column}
  \end{columns}
  \note{Slow down here even if you are behind. More than any other slide, this
  one decides whether their numbers mean anything, and you can cut the
  determinism bonus to buy the time.

  The one-off work on request one, named properly in case anyone asks: CUDA
  graph capture, kernel autotuning, lazy weight loading, and an empty prefix
  cache. The prefix cache is a block-level reuse of KV entries for a shared
  prompt prefix, which is why a repeated prompt skips prefill entirely.

  The sentence to repeat: a benchmark is a controlled experiment, and the
  control is the output length. If FP8 answers in 90 tokens and bf16 in 130,
  the end-to-end comparison is meaningless, and ignore-eos is the one-line fix.

  While circulating this is mistake number one. Look over shoulders for a
  timing loop whose first iteration is twenty times the others. If you see it,
  ask: how long did your first request take compared with your tenth, and which
  one are you reporting?}
\end{frame}

\begin{frame}[shrink=12]{Measure honestly, 2: TTFT, TPOT, concurrency}
  \centering
  \small
  \begin{tabular}{@{}L{2.4cm} L{4.1cm} L{4.3cm}@{}}
    \toprule
    \textbf{Number} & \textbf{What it is} & \textbf{Whose experience}\\
    \midrule
    \textbf{TTFT} & waiting in the queue, then reading your prompt, until the
      first token leaves & the user staring at a blank box\\
    \addlinespace[0.3em]
    \textbf{TPOT} & milliseconds per token once it is going &
      the user reading the answer as it streams\\
    \addlinespace[0.3em]
    \textbf{Throughput} & total output tokens/s across all requests at once &
      nobody; it is the \emph{operator's} number\\
    \bottomrule
  \end{tabular}
  \vspace{0.4em}

  \raggedright
  \emph{In words:} the user's total wait is the time to the first token, plus
  the time per token times the number of tokens.
  \vspace{0.3em}

  \textbf{Sweep concurrency at 1 / 8 / 32}, meaning that many requests in
  flight at the same instant. Expect throughput to rise while TPOT gets worse,
  which is Tuesday's latency--throughput curve drawn with your own numbers.
  \vspace{0.2em}

  \small
  \textbf{Sanity check at concurrency 1:} every token re-reads the whole model
  out of GPU memory, so tokens/s cannot beat about 500/s for an 8\,GB model on
  a GH200.
  \source{Definitions and the roofline estimate as given in Tuesday's lecture, section 2; the latency--throughput curve, section 4.}
  \note{Give them a feel for TPOT: 20\,ms per token is 50 tokens a second,
  comfortably faster than anyone reads.

  Say what the sweep is for. With one request you are measuring the model
  itself; at 32, most of what you see is the scheduler coping.

  The relations written symbolically, if anyone wants them on the board:
  end-to-end latency $=$ TTFT $+$ TPOT $\times$ (output tokens), and perceived
  per-user speed is $1/\mathrm{TPOT}$. The sanity check is the memory-bound
  side of the roofline from Tuesday: at batch size one, decode reads every
  weight per token, so tokens/s $\lesssim$ (memory bandwidth) $/$ (bytes of
  weights). A GH200 has roughly 4\,TB/s of HBM; an 8\,GB bf16 model therefore
  caps near 500 tokens/s, and anyone reporting more has either batched or
  mis-measured.

  Make them say which metric they would optimise for their own task. A Finnish
  NER pipeline over a document archive cares only about throughput. For a chat
  demo it is TTFT that matters. The recommendation at the end of the lab has to
  name one.

  Concurrency 32 is where the GH200 gets interesting, because the KV cache runs
  out long before the weights become a problem. Connect that back to Tuesday in
  one line and move on.

  Watch for: people running the concurrency sweep sequentially in a for loop
  and calling it concurrency 32. They need threads, or asyncio, or the vLLM
  benchmark script. Ask them how many requests the server saw at the same
  instant.}
\end{frame}

\begin{frame}{Measure honestly, 3: report the distribution}
  \begin{columns}[T]
    \begin{column}{0.52\textwidth}
      Record \key{one row per request}, not an average:
      \vspace{0.3em}
      \begin{itemize}
        \item request id, concurrency level, configuration
        \item TTFT, TPOT, output tokens, wall-clock start
      \end{itemize}
      \vspace{0.5em}
      Then report \key{median (p50) and p95} for each cell of the table, from
      \textbf{at least 50 requests}. State how many.
      \vspace{0.5em}
      The production version of this is \key{goodput}: requests per second that
      meet a latency target you choose and state, say first token under
      500\,ms \textbf{and} under 50\,ms per token after that.
      \vspace{0.3em}
      \small Those thresholds are only an example. Pick your own and justify
      them from your task.
    \end{column}
    \begin{column}{0.48\textwidth}
      \begin{gotcha}[The mean hides the tail users complain about]
        A mean latency over a batch of requests is dominated by the many fast
        ones, so it tells you nothing about the slow few.
        \vspace{0.3em}
        A server with a good mean and a terrible p95 still feels broken to use.
      \end{gotcha}
      \vspace{0.4em}
      \small
      Keep the raw per-request log in the repo. The report shows p50 and p95,
      and the log is where anyone can check them.
    \end{column}
  \end{columns}
  \note{Thirty seconds on the gotcha, and make it concrete: one request in
  twenty taking four seconds is invisible in the mean, and it is what people
  email you about.

  This is also a reproducibility point for the rubric. The raw log is what lets
  another student recompute their p95.

  Watch for: a table with one number per cell and no sample size. Ask how many
  requests that number came from. If the answer is ``one'', send them back to
  the harness. Fifty is the floor because a p95 is literally the third-worst of
  fifty; below that the statistic is noise.
  By minute 35 the bf16 row should be complete.}
\end{frame}

% =====================================================================
\section{Quality}
% =====================================================================

% The exact server command is maintained in lab3-example/scripts.
\iffalse
\begin{frame}[fragile]{FP8 on the GH200: one flag}
\begin{lstlisting}[language=slurm]
# Stop the bf16 server first (Ctrl-C). The port and the GPU memory are not shared.
# --quantization fp8   store numbers in 8 bits instead of 16. Native on the GH200.
# --kv-cache-dtype fp8 optional, and a SEPARATE experiment -- one change at a time

vllm serve Qwen/Qwen3.5-4B-Instruct \
  --served-model-name base \
  --enable-lora \
  --lora-modules task=/scratch/project_2021893/lora/w2-adapter \
  --quantization fp8 \
  --max-model-len 4096 \
  --gpu-memory-utilization 0.85 \
  --port 8000
\end{lstlisting}
  \vspace{0.2em}
  \small
  Then re-run the \emph{identical} timing cells. Do not edit the harness
  between rows.
  \vspace{0.3em}

  What you should see, and should verify rather than assume: an 8\,GB bf16
  model drops to about 4\,GB, the freed 4\,GB becomes KV cache, so concurrency
  32 improves more than concurrency 1 does.
  \source{FP8 is native on H100-class hardware (GH200); see Tuesday's lecture, section 5.
  The direction is what you should expect. How big it turns out to be is your result.}
  \note{FP8 is a one-flag change on this hardware, which is exactly why it is
  the right experiment for a 90-minute lab. The format is E4M3 --- four
  exponent bits, three mantissa bits --- applied to weights and activations;
  the KV cache stays 16-bit unless they also pass \texttt{--kv-cache-dtype}.

  The point to make about the expected result: decode is bandwidth-bound, so
  halving the bytes read per token is the whole mechanism. Anybody who sees no
  speed-up at concurrency 1 should check they are not prompt-bound with a
  128-token output.

  Say clearly: kv-cache-dtype fp8 is optional and is a separate experiment. If
  they turn it on in the same run they have changed two things again.

  Minute 35 to 60. If a student is still fighting the bf16 row at minute 45,
  tell them to skip the concurrency sweep and get the FP8 row instead. Two
  filled rows are worth more than three concurrency levels in one row.}
\end{frame}

\fi
\begin{frame}[shrink=14]{Score the quantised model yourself}
  \begin{columns}[T]
    \begin{column}{0.54\textwidth}
      Run the \key{same held-out scorer used in week 2} against the FP8
      endpoint. It should take a model or a base URL and return the task score
      in one call.
      \vspace{0.4em}

      Same held-out set, same prompts, same decoding parameters as the bf16
      score. Otherwise you are comparing two different experiments.
      \vspace{0.6em}

      \textbf{Where quantisation damage actually shows up}
      \begin{itemize}
        \item long inputs
        \item non-English text
        \item long reasoning chains
      \end{itemize}
      \vspace{0.2em}
      \small
      Headline benchmarks are the \emph{last} place it appears, which is why a
      published ``FP8 is lossless'' claim says nothing about your task.
    \end{column}
    \begin{column}{0.46\textwidth}
      \begin{block}{Report both halves}
        \small
        \textbf{Gained:} tokens/s, p95 latency, memory freed for KV cache,
        concurrency reached.\\[0.4em]
        \textbf{Paid:} score points on the held-out set.\\[0.4em]
        If you report only one of these halves, you have written an advert.
      \end{block}
      \vspace{0.4em}
      \small
      Look at your \textbf{ten worst examples} under FP8, not only the
      aggregate. If the aggregate moved by less than the noise of your metric,
      say so.
    \end{column}
  \end{columns}
  \source{Quality cost of quantisation as discussed in Tuesday's lecture, section 5.
  The score change on your task is a measurement, not a citation.}
  \note{This is the pedagogical centre of the second half. The transferable
  habit is: when a vendor says a compression method is free, score it on your
  own data, on your own hardware, with your own metric.

  The non-English point lands with the Finnish projects in the room. Ask one of
  the TyDi QA or Turku NER students to say their FP8 delta out loud when they
  have it.

  Watch for: students scoring FP8 on a different number of examples than they
  scored bf16 on, usually because they truncated the set to save time. Either
  truncate both or neither.}
\end{frame}

% The schema and request implementation are in lab3-example.
\iffalse
\begin{frame}[fragile,shrink=5]{Constrained output: a schema and a request}
\begin{lstlisting}
schema = {
  "type": "object",
  "properties": {
    "label":      {"type": "string",
                   "enum": ["entailment", "neutral", "contradiction"]},
    "confidence": {"type": "number", "minimum": 0, "maximum": 1}
  },
  "required": ["label", "confidence"],
  "additionalProperties": False
}

r = client.chat.completions.create(
    model="task",
    messages=[{"role": "user", "content": prompt}],
    temperature=0.0, max_tokens=64,
    extra_body={"guided_json": schema},
)
\end{lstlisting}
  \small
  \emph{In words:} the server turns the schema into a rule about which tokens
  may come next, and at each step it blocks the ones that would break it.
  Nothing has to be cleaned up afterwards, because invalid JSON cannot be
  produced in the first place.
  \textbf{Success looks like} every output surviving \texttt{json.loads}.
  \note{The mechanism properly: the schema is compiled to a finite-state
  machine, and at every decoding step the logits of tokens that would leave the
  accepting path are set to minus infinity before the softmax. That is why the
  invalid rate goes to essentially zero rather than merely down.

  Replace the enum with their own label set: NLI labels, NER types, an answer
  span plus a confidence, whatever their task needs. The schema should be the
  smallest thing that captures their output format.

  Watch for: schemas with no enum and no required list, which constrain almost
  nothing and then surprise the student when the invalid rate barely moves.
  Also watch for max-tokens too small, which truncates valid JSON and shows up
  as an invalid output.}
\end{frame}

\fi
\begin{frame}{Measure the invalid-output rate, both ways}
  \begin{columns}[T]
    \begin{column}{0.52\textwidth}
      \textbf{The protocol}
      \begin{enumerate}
        \item A fixed set of held-out prompts, asking for JSON in the prompt
          text, \key{no} schema. Parse each output. Count failures.
        \item The \emph{same} prompts with \texttt{guided\_json}. Count again.
        \item Score \emph{content} in both runs with the same held-out scorer.
      \end{enumerate}
      \vspace{0.5em}
      Invalid means: not parseable, or parseable but missing a required field,
      or a label outside your enum. Say which definition you used.
    \end{column}
    \begin{column}{0.48\textwidth}
      \begin{block}{Record two numbers}
        \small
        \textbf{Syntax:} invalid-output rate, with and without the schema.\\[0.4em]
        \textbf{Content:} task score, with and without the schema.
      \end{block}
      \vspace{0.4em}
      \small
      Forcing valid syntax can cost content quality: the model gets pushed onto
      token paths it would not have chosen. If your invalid rate goes to zero
      and your task score drops two points, \textbf{that is the finding}, and
      it belongs in the recommendation.
      \vspace{0.3em}

      Constrained decoding also costs a little speed. Time it if you have the
      minutes; the score matters more.
    \end{column}
  \end{columns}
  \note{Minute 60 to 80. By now the table has two full rows; this is the third.

  The counter-intuitive bit is worth saying plainly: valid JSON with the wrong
  label inside it is still a wrong answer. The students who find a
  syntax-versus-content trade-off will have something genuinely interesting to
  write about, so encourage them to report it rather than hide it.

  Watch for: a baseline run with no instruction to produce JSON at all, which
  makes the unconstrained invalid rate near 100 per cent and the comparison
  trivial. The honest baseline is a good prompt asking for JSON.}
\end{frame}

\begin{frame}{Bonus: temperature 0, alone and in a crowd}
  \begin{columns}[T]
    \begin{column}{0.54\textwidth}
      \textbf{The procedure}
      \begin{enumerate}
        \item Pick one prompt. \texttt{temperature=0}, \texttt{max\_tokens=256}.
        \item Send it \key{alone} on an idle server. Save the output.
        \item Send it again alone. Confirm it is identical.
        \item Send the \emph{same} prompt as one of \key{64 concurrent}
          requests. Save that output.
        \item Diff against the solo output.
        \item Repeat steps 4--5 twenty times. Report how many of the 20 differ.
      \end{enumerate}
    \end{column}
    \begin{column}{0.46\textwidth}
      \begin{block}{Why it happens}
        \small
        In floating point, $(a+b)+c$ and $a+(b+c)$ can differ in the last bits,
        and the GPU adds numbers up in a different order depending on how many
        requests it is batching.
        \vspace{0.3em}
        So \emph{other people's requests} can change your scores, and
        occasionally the words you get back.
      \end{block}
      \vspace{0.3em}
      \small
      The fix is called \key{batch invariance}. This is not GPU randomness and
      it is not a seed problem; see Tuesday, section 8.
      \vspace{0.3em}

      One line in the report: what this means for reproducible evaluation.
    \end{column}
  \end{columns}
  \note{This is the bonus, so it is the first thing to cut if the room is
  behind. Say that when you introduce it, so nobody starts here.

  Precisely: the reduction order inside the attention and matmul kernels
  depends on batch shape and on the scheduler's split of work, so the logits
  differ in the last bits; once an argmax flips, the two continuations diverge
  completely. The fix is batch-invariant kernels rather than a seed.

  The reason it is worth doing at all: they have spent three weeks reporting
  scores to two decimal places, and this shows them that the third decimal
  depends on who else was on the server. That reframing is the point.

  If someone sees zero differences in twenty trials, that is a real result
  worth a sentence: their outputs may be short, highly constrained, or the
  divergence point is past 256 tokens. Ask them to try a longer generation.}
\end{frame}

% =====================================================================
\section{Admin}
% =====================================================================

% The restart scripts are maintained in lab3-example/scripts.
\iffalse
\begin{frame}[fragile,shrink=5]{A server you can restart with one command}
\begin{lstlisting}[language=slurm]
#!/bin/bash
# scripts/serve.sh --- must work from a cold start, on a clean shell
set -euo pipefail                  # stop at the first error, not three steps later

module purge                       # do not inherit whatever was loaded before
module load python-vllm

MODEL=${MODEL:-Qwen/Qwen3.5-4B-Instruct}   # ${X:-default}: use $X if set, else this
ADAPTER=${ADAPTER:-$PWD/lora/w2-adapter}
PORT=${PORT:-8000}

vllm serve "$MODEL" \
  --served-model-name base \
  --enable-lora --lora-modules task="$ADAPTER" \
  --max-model-len 4096 \
  --gpu-memory-utilization 0.85 \
  --port "$PORT" "$@"              # "$@" passes any extra flags straight through
\end{lstlisting}
  \small
  \key{Test it properly:} kill your server, open a \emph{fresh} terminal, run
  \texttt{bash scripts/serve.sh}, and confirm that
  \texttt{curl localhost:8000/v1/models} lists \texttt{base} and
  \texttt{task}. No manual steps first, nothing recalled from your history.
  \note{Be firm about this one. In week 5 Dmitry's RAG lab uses this endpoint
  as its generator, and in week 6 the agent calls it. A script that only works
  because of something the student typed earlier will cost them a lab.

  The cold-start test is the whole instruction: a fresh terminal, no exports,
  and no accidental \texttt{cd} into the right directory. It is ninety seconds
  of work.

  Also tell them to put the base URL and the model names in the repo README.
  Dmitry's notebooks will ask for them.}
\end{frame}

\fi
\begin{frame}{Files and results to keep}
  \begin{columns}[T]
    \begin{column}{0.52\textwidth}
      \textbf{In the repo}
      \begin{itemize}
        \item \texttt{results/speed-quality.csv}, the table with one row per
          configuration
        \item \texttt{results/requests.csv}, the raw per-request log
        \item \texttt{scripts/serve.sh}, tested from a cold start
        \item the notebook with the timing harness
        \item a \texttt{README} giving the base URL, port, model and adapter
          names
      \end{itemize}
      \vspace{0.5em}
      \textbf{In the report, section 4}\\
      The table, and 3--5 sentences of recommendation.
    \end{column}
    \begin{column}{0.48\textwidth}
      \begin{block}{The recommendation must answer}
        \small
        Which configuration would you deploy, and under what latency target?\\[0.3em]
        What did FP8 cost you, in score points?\\[0.3em]
        What would you change at ten times the traffic?
      \end{block}
      \vspace{0.4em}
      \small
      There is no intermediate submission. Keep these files for the final
      repository and reuse the endpoint in later labs if it suits your project.
      \vspace{0.3em}

      Rubric reminder: Results 30\,\%, Analysis 30\,\%. The table counts as
      Results and the recommendation as Analysis.
    \end{column}
  \end{columns}
  \note{Two minutes at most, and it can be the last thing you say before
  releasing the room if you are short of time.

  The three questions in the block are deliberately the same three questions a
  colleague would ask in a code review. Students who can answer them have
  understood the lab.

  Emphasise that saving an exact server command makes later experiments much
  easier to reproduce.}
\end{frame}

\begin{frame}{GPU etiquette}
  \begin{columns}[T]
    \begin{column}{0.52\textwidth}
      \begin{itemize}
        \item \key{Never serve from a login node.} No GPU there, and the CPUs
          are shared by everyone on Roihu.
        \item \texttt{gputest} (15 min) before \texttt{gpumedium} (36 h) for
          anything you submit as a batch job.
        \item A \texttt{gpuinteractive} session holds a whole GH200 for up to
          12 hours. \key{Close it when you leave today}, from the Roihu web
          interface rather than by closing the browser tab.
        \item Stop your vLLM server before you stop the session, so the port
          and the memory are released cleanly.
      \end{itemize}
    \end{column}
    \begin{column}{0.48\textwidth}
      \begin{block}{Why it matters more this week}
        \small
        We share one course project allocation, and \textbf{weeks 5 and 6 need
        this endpoint again}. Every forgotten session today is GPU time missing
        from Dmitry's labs.
      \end{block}
      \vspace{0.5em}
      \small
      Your server dies with the allocation. That is expected, and it is why
      \texttt{serve.sh} has to work from a cold start.
      \vspace{0.3em}

      \texttt{squeue --me} to see what you are holding.
      \texttt{scancel <jobid>} to let it go.
    \end{column}
  \end{columns}
  \note{Short and firm, thirty seconds. The line that changes behaviour is the
  connection to weeks 5 and 6: the GPU time wasted today is their own compute
  later in the term.

  At the end of the session, before people drift off, walk the room once and
  ask each table whether their session is closed. Do it out loud so the rest of
  the room hears.

  Mention that we can all see each other's jobs in the queue. That tends to be
  more effective than the policy itself.}
\end{frame}

\begin{frame}{On the wall for the rest of the session}
  \begin{columns}[T]
    \begin{column}{0.45\textwidth}
      \textbf{Timeline}
      \vspace{0.2em}

      \small
      \begin{tabular}{@{}r L{3.4cm}@{}}
        \toprule
        \textbf{min} & \\
        \midrule
        0--10  & server up, adapter visible\\
        10--35 & speed: 1 / 8 / 32\\
        35--60 & FP8: re-time, re-score\\
        60--80 & JSON schema: invalid rate\\
        80--90 & save results and write up\\
        \bottomrule
      \end{tabular}
      \vspace{0.6em}

      \footnotesize
      Server not up by minute 15? Hand up.\\
      Bonus (determinism) only if rows 1--3 are done.
    \end{column}
    \begin{column}{0.55\textwidth}
      \textbf{Done when}
      \vspace{0.2em}

      \small
      \texttt{[\,]}\ server answers on \texttt{:8000}, model \texttt{task}\\
      \texttt{[\,]}\ warm-up requests discarded\\
      \texttt{[\,]}\ p50 \emph{and} p95 for TTFT and TPOT, at 1 / 8 / 32\\
      \texttt{[\,]}\ number of requests recorded for every cell\\
      \texttt{[\,]}\ FP8 row: timed \emph{and} scored\\
      \texttt{[\,]}\ invalid-output rate, with and without schema\\
      \texttt{[\,]}\ \texttt{results/speed-quality.csv} saved\\
      \texttt{[\,]}\ \texttt{scripts/serve.sh} tested from a clean shell\\
      \texttt{[\,]}\ 3--5 sentence deployment recommendation\\
      \texttt{[\,]}\ GPU session closed
      \vspace{0.5em}

      \footnotesize Keep these outputs for the final report and repository.
    \end{column}
  \end{columns}
  \note{This slide stays projected for the whole working part of the session.
  Do not advance past it. It is the operating manual: students will look up at
  it instead of asking where they should be.

  Say one sentence when you put it up --- ``this is on the wall, check yourself
  against it'' --- then stop talking and start circulating. The briefing ends
  here, at minute 10 to 15.

  Use the checklist as your own circulating script too. Walking up to a table
  and asking ``which of these ten is your next one'' is a better opening than
  ``how is it going'', and it surfaces the stuck students faster.}
\end{frame}

\appendix

\begin{frame}{When it breaks}
  \scriptsize
  \begin{tabular}{@{}L{3.4cm} L{7.6cm}@{}}
    \toprule
    \textbf{Symptom} & \textbf{First thing to try}\\
    \midrule
    OOM when the server starts &
      Something else holds the GPU, usually a week 2 notebook kernel.
      Restart the kernel, then run \texttt{nvidia-smi}, which lists what is
      holding the card. Then lower \texttt{--gpu-memory-utilization} to 0.7 and
      \texttt{--max-model-len} to 2048.\\
    \addlinespace[0.3em]
    \texttt{Address already in use} &
      Your previous server is still running. \texttt{pkill -f "vllm serve"}, or
      use \texttt{--port 8001} and change the base URL to match.\\
    \addlinespace[0.3em]
    404, or the adapter has no effect &
      \texttt{model=} must be the name on the \emph{left} of
      \texttt{--lora-modules name=path}, not the repo id. Check
      \texttt{--enable-lora} is present and the path contains
      \texttt{adapter\_config.json}.\\
    \addlinespace[0.3em]
    Fluent nonsense, normal latency &
      Tokenizer or chat-template mismatch: the adapter was trained on the
      instruct model and you are serving the base one, or the reverse. Serve
      the checkpoint you fine-tuned.\\
    \addlinespace[0.3em]
    Server dies when the allocation ends &
      Expected: the server is a child of your
      \texttt{gpuinteractive} session. Start a new session and run
      \texttt{bash scripts/serve.sh}.\\
    \bottomrule
  \end{tabular}
  \note{Keep this slide one key-press away all session. Five symptoms cover
  almost everything that goes wrong in this lab.

  The fourth row is the subtle one and the easiest to miss over a shoulder,
  because the latency numbers look perfectly healthy. If a student's speed
  table is fine but their score collapsed against the bf16 baseline, check the
  template before anything else.

  The fifth row is worth announcing to the whole room once, around minute 70,
  because the reservation window ends and several servers will die at the same
  time. It also makes the argument for the restart script for you.}
\end{frame}

\end{document}
